\documentclass[fontsize=12pt, paper=a4]{scrartcl}
\usepackage{amssymb, amsthm, amsmath, graphicx, lmodern, url, orcidlink}
\usepackage[margin=2.5cm]{geometry}

\begin{document}

\title{Gravity and Action from the Evolution of a Complex Manifold}
\author{Donald Airey~\orcidlink{0000-0002-2958-1545}}
\date{}

\maketitle
\begin{abstract}
We present a bilocal geometric framework in which separations are defined between ordered pairs of events rather than by a local infinitesimal metric. This construction is motivated by the need to represent global evolution in a manner that remains well-defined over finite separations while admitting a smooth local limit. The resulting geometry is formulated on a complex manifold with an imaginary temporal coordinate and a real spatial coordinate. A constant manifold acceleration sets the global spatial scale and produces a parabolic cycle of expansion and contraction. In the local reduction, the same acceleration scale yields a circular-orbit envelope consistent with the baryonic Tully--Fisher relation. At cosmological scales, the resulting analytic distance--redshift relation predicts the Pantheon+ Type~Ia supernova luminosity distances on the SH0ES-calibrated absolute scale and yields a lower covariance-weighted $\chi^2$ than Planck-calibrated FLRW, without invoking dark matter or dark energy. The parabolic metric evolution, when combined with standard recombination microphysics, yields an acoustic angular scale consistent with Planck measurements.
\end{abstract}

\begin{figure}
\centering
\includegraphics[width=0.6\linewidth]{Manifold.pdf}
\caption{Schematic representation of the complex parabolic metric manifold. Latitude lines correspond to spatial slices at fixed epoch, while longitude lines trace the temporal evolution of the manifold.}
\label{fig:manifold}
\end{figure}

\section{Introduction}

Cosmological observations tightly constrain theoretical descriptions of the universe. Measurements of the cosmic microwave background (CMB) fix early-universe conditions and global geometry with high precision~\cite{Planck2018}, while Type~Ia supernovae (SNe~Ia) probe the late-time expansion history through the luminosity distance--redshift relation~\cite{Brout2022}. Although each dataset is internally consistent, discrepancies arise when parameters inferred from early- and late-time observations are compared. In particular, the tension in the Hubble constant~\cite{Planck2018,Riess2022} indicates that the standard cosmological model may be incomplete.

Historically, such discrepancies have been addressed by extending the model. The cosmological constant was introduced to permit a static universe~\cite{Einstein1917} and abandoned following the discovery of expansion~\cite{Hubble1929}. Galaxy rotation curves motivated the introduction of non-baryonic dark matter~\cite{Rubin1980}, and supernova evidence for accelerated expansion led to the reintroduction of the cosmological constant as dark energy~\cite{Riess1998,Perlmutter1999}. While these additions yield a phenomenologically successful framework, the dominant components of the cosmic energy budget remain physically unexplained~\cite{Planck2018}. More recent approaches introduce additional degrees of freedom, including time-varying dark energy~\cite{DiValentino2021}, further increasing the parametric structure of the model.

We develop an alternative framework, termed Parabolic Metric Evolution (PME), in which large-scale evolution and local gravitational phenomena arise from a common geometric structure governed by a single intrinsic acceleration scale. The framework rests on two structural choices. First, separations are defined bilocally between ordered pairs of events rather than through a local infinitesimal metric, so that finite intervals remain well-defined even when the geometry evolves on the scale of the separation itself. Second, the manifold evolution is governed by a constant homogeneous acceleration $A$ --- a single intrinsic scale, the minimum required for nontrivial dynamics, whose value is determined empirically. The resulting homogeneous evolution yields a parabolic manifold extent, motivating the name Parabolic Metric Evolution (Figure~\ref{fig:manifold}).

These two choices fix the framework's predictions. The constant acceleration scale yields a circular-orbit envelope that reproduces the baryonic Tully--Fisher relation, with $A$ determined empirically from galaxy rotation curves. Applied to the analytic distance--redshift relation that follows from the bilocal interval, the same scale predicts the Pantheon+ Type~Ia supernova luminosity distances on the SH0ES-calibrated absolute scale, achieving a lower covariance-weighted $\chi^2$ than Planck-calibrated FLRW. The present-day Hubble constant emerges as a prediction rather than a fitted parameter, and the acoustic angular scale at recombination is reproduced when the manifold kinematics are combined with standard recombination microphysics. No independent dark matter or dark energy sectors are required; the same acceleration scale governs galactic dynamics, cosmological expansion, and the baryonic density.

The framework differs from relativistic metric gravity in its bilocal definition of separation, and from modified-dynamics frameworks such as MOND~\cite{Milgrom1983,McGaugh2016} in that the characteristic acceleration scale is not introduced phenomenologically but inherited from the global manifold kinematics. Cosmological evolution, gravity, and local dynamics are thereby unified through a single geometric scale.

\section{Bilocal Geometry}

\subsection{Foundational Structure}

The framework is defined by the following postulates:

\begin{enumerate}
\item Physical separations are defined bilocally between ordered pairs of events.
\item The manifold evolution is governed by a constant acceleration scale $A$.
\item The manifold is parameterized by an intrinsic evolution coordinate $\chi$.
\item Observable time is defined operationally through null exchange, providing an invariant ordering of events. The intrinsic evolution parameter is related to this observable time by the projection $\chi = i t$, which fixes the real--imaginary decomposition of temporal and spatial functionals and serves as a defining structural postulate of the framework.
\end{enumerate}

Within this framework, geometric separation is defined operationally through ordered bilocal relations and possesses no independently specified continuum structure beneath this construction.  All subsequent results follow from this structure.

\subsection{Bilocal Spatial Separation}

Let $p_e=(x_e^0,x_e^1)$ denote an emission event and $p_o=(x_o^0,x_o^1)$ an observation event, with the bilocal interval constructed from a temporal functional obtained from accumulated evolution between slices and a spatial functional defined from separations on the emission and observation slices.

The bilocal construction is illustrated in Figure~\ref{fig:bilocalconstruction}. The spatial slices are closed, with $L(t)$ denoting the circumference of the slice at epoch $t$. The spatial slices scale with this global manifold extent, so the separations at emission and observation satisfy
\begin{equation}
\frac{\Delta x_e^1}{L(t_e)}=\frac{\Delta x_o^1}{L(t_o)} .
\label{eq:comoving_definition}
\end{equation}
Imposing symmetry under interchange of emission and observation together with a smooth coincidence limit motivates adopting the midpoint average as the minimal symmetric choice,
\begin{equation}
\Delta s^1=\Delta x_o^1-\frac12(\Delta x_o^1-\Delta x_e^1),
\end{equation}
\begin{equation}
\Delta s^1=\frac12(\Delta x_o^1+\Delta x_e^1),
\label{eq:spatial_half_expansion}
\end{equation}

\begin{figure}
\centering
\includegraphics[width=0.7\linewidth]{BilocalGeometry.pdf}
\caption{Geometric construction of the bilocal interval between emission event $p_e$ and observation event $p_o$, showing the accumulated temporal evolution and the symmetric midpoint spatial separation between scaled slices.}
\label{fig:bilocalconstruction}
\end{figure}

\subsection{Invariance of the Bilocal Interval}
The bilocal interval must be a scalar depending only on the functionals $(\Delta s^0,\Delta s^1)$. Four constraints uniquely fix its form. First, minimality: the lowest-order scalar is the general quadratic
\begin{equation}
\Delta s^2 = a\,(\Delta s^0)^2 + b\,(\Delta s^1)^2 + 2c\,\Delta s^0 \Delta s^1 ,
\end{equation}
with higher-order terms excluded. Second, non-degeneracy: the null condition must yield a finite null ratio $\Delta s^1/\Delta s^0$ in the local limit, requiring that neither term vanish from the interval, so $a\neq 0$ and $b\neq 0$. Third, endpoint-interchange symmetry: under interchange of $(p_e,p_o)$ the temporal functional changes sign while the symmetric spatial functional does not, so the mixed term changes sign and invariance requires $c=0$. Fourth, canonical normalization: $\Delta s^0$ and $\Delta s^1$ are identified with the imaginary and real components of the manifold respectively, and the complex structure fixes their relative weighting to unity. These four constraints together select the unique bilocal interval
\begin{equation}
\Delta s^2 = (\Delta s^0)^2 + (\Delta s^1)^2 ,
\label{eq:bilocal_interval}
\end{equation}
with $\Delta s^0$ purely imaginary and $\Delta s^1$ real.

\subsection{Intrinsic Tangent Evolution}

The intrinsic evolution of the manifold is parameterized by a coordinate $\chi$. Observable time is introduced by the projection $\chi = i t$, so that $d\chi = i\,dt$, fixing the real--imaginary decomposition used in the bilocal construction. With constant intrinsic acceleration $A$, the tangent evolution scale along the manifold history follows from integration with respect to $\chi$,
\[
v_\chi(t) = \int A\, d\chi .
\]
Substituting $d\chi = i\,dt$ gives $v_\chi(t) = i\!\int A\,dt$, which evaluates to
\begin{equation}
v_\chi(t) = i(At - V_0),
\label{eq:imaginary_velocity}
\end{equation}
where $V_0$ is an integration constant representing the initial expansion rate of the manifold. The intrinsic tangent evolution is therefore purely imaginary, in accordance with the real--imaginary decomposition fixed by the projection $\chi = i t$. Because the acceleration $A$ acts homogeneously across the manifold, this tangent evolution is universal: every point on the manifold at epoch $t$ shares the same intrinsic tangent generator $i(At - V_0)$. This tangent structure is physically real but not directly observable.

\subsection{Induced Spatial Evolution}

The observable spatial extent arises from integrating the intrinsic tangent evolution along the imaginary temporal direction. The global spatial scale of the manifold is therefore $L(t) \equiv \int v_\chi(t)\, d\chi$. Substituting the expression for $v_\chi(t)$ and using $d\chi = i\,dt$ gives $L(t) = \int i(At - V_0)\, i\, dt$. Evaluating the integral yields
\begin{equation}
L(t) = V_0 t - \frac{1}{2} A t^2.
\label{eq:manifold_extent}
\end{equation}

\subsection{Manifold Deceleration}

The intrinsic evolution is defined along $\chi$, while observable dynamics are expressed in the projected time coordinate $t$ with $d\chi = i\,dt$. Rewriting the manifold extent in terms of $t$ yields its observable form. Differentiating twice with respect to $t$ gives
\begin{equation}
\frac{d^2 L}{dt^2} = -A,
\label{eq:manifolddeceleration}
\end{equation}
so the manifold extent evolves with a constant kinematic deceleration of magnitude $A$, which is universal and independent of position or epoch.

\section{Cosmological Distance}

Observable cosmological distances follow from the bilocal interval defined in equation~(\ref{eq:bilocal_interval}). 

\subsection{Temporal Separation}

The temporal component of the bilocal interval is obtained from the accumulated tangent evolution along the manifold history between the emission and observation parameters $t_e$ and $t_o$, $\Delta s^0 \equiv \int_{t_e}^{t_o} v_\chi(t)\,dt$. Substituting Eq.~(\ref{eq:imaginary_velocity}) gives
\begin{equation}
\Delta s^0 = \int_{t_e}^{t_o} i\left(A t - V_0\right)\,dt .
\end{equation}
Evaluating the integral yields
\begin{equation}
\Delta s^0 = i\left(\frac{A}{2}\left(t_o^2 - t_e^2\right) - V_0\left(t_o - t_e\right)\right) .
\end{equation}
The temporal separation is purely imaginary and depends only on the endpoints.

\subsection{Spatial Separation}

Solving equation~(\ref{eq:comoving_definition}) for the emission separation gives
\begin{equation}
\Delta x_e^1 = \Delta x_o^1\,\frac{L(t_e)}{L(t_o)} .
\end{equation}

Because spatial slices scale with the manifold extent $L(t)$ and the slice separations satisfy the bilocal scaling relation given in equation~(\ref{eq:spatial_half_expansion}), the spatial component becomes
\begin{equation}
\Delta s^1 =\frac{\Delta x_o^1}{2}\left(1 +\frac{V_0 t_e - \frac{1}{2} A t_e^2}{V_0 t_o - \frac{1}{2} A t_o^2}\right).
\label{eq:bilocal_spatial}
\end{equation}

\subsection{Bilocal Metric Formula}

Substituting the temporal and spatial separations into equation~(\ref{eq:bilocal_interval}) yields the interval relating the emission and observation events,
\begin{equation}
\Delta s^2 = \left[i\left(\frac{A}{2}\left(t_o^2 - t_e^2\right) - V_0\left(t_o - t_e\right)\right)\right]^2+ \left[\frac{\Delta x_o^1}{2}\left(1 + \frac{V_0 t_e - \frac{1}{2}A t_e^2}{V_0 t_o - \frac{1}{2}A t_o^2}\right)\right]^2 .
\label{eq:bilocal_metric}
\end{equation}

\subsection{Redshift Relation}

Redshift is defined observationally by the ratio of observed to emitted wavelength,
\begin{equation}
1+z \equiv \frac{\lambda_o}{\lambda_e}.
\end{equation}
We identify wavelength with a comoving spatial separation and assume that such separations scale with the manifold extent, so that $\lambda(t) \propto L(t)$.
\begin{equation}
1+z = \frac{L(t_o)}{L(t_e)}.
\label{eq:redshift_relation}
\end{equation}

Using equation~(\ref{eq:manifold_extent}), the emission epoch $t_e$ is expressed in terms of $t_o$ and the observed redshift $z$,
\begin{equation}
t_e = \frac{V_0(1+z) - \sqrt{(1+z)\big[(V_0 - A t_o)^2 + V_0^2 z\big]}}{A(1+z)},
\label{eq:te_solution}
\end{equation}
where the negative branch is chosen so that $t_e < t_o$.

Substituting into equation~(\ref{eq:bilocal_metric}) and imposing the null condition $\Delta s^2 = 0$ yields the observable spatial separation,
\begin{equation}
D_C(z)= \Delta x_o^1= \frac{t_o\left(2V_0 - A t_o\right) z}{2 + z}.
\label{eq:distance_redshift_relation}
\end{equation}

\section{Causal Structure}

The null structure of the bilocal interval also determines the extent of causal connectivity on the manifold. The largest spatial separation permitted by null ordering is obtained from equation~(\ref{eq:bilocal_metric}) with the emission event at the origin of the manifold history, $t_e = 0$, and the observation event at epoch $t_o = t$. Solving the null condition $\Delta s^2 = 0$ for the observed separation gives the causal horizon
\begin{equation}
d_{\mathrm{ch}}(t) = t\left(2V_0 - A t\right) .
\end{equation}

Differentiating the causal horizon gives the expansion rate of the causal boundary,
\begin{equation}
c_{\mathrm{c}}(t)=\frac{d}{dt}\,d_{\mathrm{ch}}=2\left(V_0-At\right).
\end{equation}
This quantity defines the rate at which the accessible causal domain can be extended under the operational clock, and therefore characterizes the growth of the causally orderable region.

Comparing the causal horizon with the manifold extent $L(t)$ defined in equation~(\ref{eq:manifold_extent}) yields
\begin{equation}
\frac{d_{\mathrm{ch}}(t)}{L(t)} = 2 .
\end{equation}
The causal horizon is therefore twice the manifold extent at all epochs. This ratio is constant and position-independent, since $A$ acts homogeneously across every spatial slice: every point at every epoch shares the same causal relationship to the rest of the manifold. Large-scale homogeneity and isotropy of the cosmic microwave background follow directly from the geometry.

\section{Gravity}

When the temporal interval between neighboring events is small compared with the characteristic evolution scale of the manifold extent $L(t)$, the bilocal construction admits a smooth local reduction yielding an effective acceleration field inherited from the manifold kinematics.

\subsection{Coincidence Limit}

The local description is obtained as the coincidence limit of the bilocal interval as the endpoint separation tends to zero. Let the observation event occur at epoch $t$ and the emission event at $t-\Delta t$, with $\Delta t \to 0$. To extract the leading nontrivial local structure, the temporal and spatial bilocal functionals are expanded to first and second orders in $\Delta t$.

Using the local expansion of the manifold extent,
\[
L(t-\Delta t)=L(t)-\dot L(t)\,\Delta t-\frac{1}{2}A\,\Delta t^2,
\]
the slice-separation relation gives
\begin{equation}
\frac{\Delta x_e^1}{L(t-\Delta t)}=\frac{\Delta x_o^1}{L(t)} .
\end{equation}
Hence
\begin{equation}
\Delta x_e^1=\Delta x_o^1\left(1-\frac{\dot L}{L}\Delta t\right),
\end{equation}
so the bilocal spatial component becomes
\begin{equation}
\Delta s^1=\Delta x_o^1\left(1-\frac{1}{2}\frac{\dot L}{L}\Delta t\right)\approx \Delta x_o^1
\end{equation}
to leading order in $\Delta t$. In this limit $\Delta s^1$ coincides with the observed separation, and we write $\Delta x \equiv \Delta x_o^1$ for the local observable displacement. Over the same interval, the temporal component is
\begin{equation}
\Delta s^0=\int_{t-\Delta t}^{t} i\left(A t' - V_0\right)\,dt' \approx i\left(At-V_0\right)\Delta t.
\end{equation}
Substituting these leading terms into equation~(\ref{eq:bilocal_interval}) yields
\begin{equation}
\Delta s^2 \approx -\left(At-V_0\right)^2\Delta t^2 + \Delta x^2 ,
\label{eq:local_interval}
\end{equation}
which is the emergent local interval obtained from the coincidence limit of the bilocal geometry. Although the intrinsic temporal contribution is imaginary, the squared interval is real, and the observable null scale arises from the operational projection condition that defines admissible signal exchange. Imposing $\Delta s^2 = 0$ on the coincidence-limit interval gives
\begin{equation}
\left(At-V_0\right)^2 \Delta t^2 = \Delta x^2 ,
\end{equation}
so the null ratio yields a real observable scale,
\begin{equation}
\frac{\Delta x}{\Delta t} = V_0 - At ,
\label{eq:real_tangent_scale}
\end{equation}
where the positive root is selected by the requirement that signal propagation proceed in the direction of increasing time on the expansion branch $At < V_0$; the contraction branch lies beyond the scope of the present treatment. The quantity $(V_0-At)$ is therefore the observable null scale at epoch $t$, emerging as the real operational projection of the underlying imaginary tangent generator through the null condition that defines admissible exchange.

\subsection{Acceleration Field}

Consider an event pair $(p_e,p_o)$ whose temporal separation $\Delta t$ satisfies $\Delta t \ll L/\dot L$. Expanding the manifold extent about a reference epoch $t_0$ gives
\begin{equation}
L(t_0+\delta t)=L(t_0)+\dot{L}(t_0)\,\delta t+\frac{1}{2}\ddot{L}(t_0)\,\delta t^2+\cdots .
\end{equation}
Using equation~(\ref{eq:manifolddeceleration}), the quadratic term is governed by the constant local value $\ddot{L}(t_0)=-A$. The local reduction therefore contains a homogeneous second-order contribution corresponding to a uniform deceleration of magnitude $A$ inherited from the global manifold evolution, which provides the physical origin of gravity in the local limit. This uniform deceleration defines the background acceleration field underlying local gravitational dynamics.

The effective inward acceleration in a bound system decomposes into a homogeneous background and a sourced concentration,
\begin{equation}
a_{\mathrm{in}}(x)=a_{\mathrm{bg}}+a_{\mathrm{conc}}(x).
\label{eq:ain_scalar}
\end{equation}
The background term $a_{\mathrm{bg}}$ is uniform and has constant magnitude $A$. In the absence of localized content, it defines the manifold evolution, and an inertial frame follows it without inducing concentration or depletion. Localized inertial content, however, resists this acceleration, redistributing the flux into spatial concentrations that define the sourced gravitational component $a_{\mathrm{conc}}(x)$ governing bound dynamics. The inertial parameter $m$ entering this redistribution is the same quantity that weights the invariant separation in the construction of geometric action, so that the gravitational response corresponds to a redistribution of the momentum-weighted action-space separation underlying the local geometry.

\subsection{Minimal Local Closure}

The bilocal construction fixes the homogeneous background acceleration scale $A$ through the second-order local reduction of the manifold extent but does not determine the sourced concentration arising from spatial variations in inertial resistance. A weak-field closure is therefore required.

We adopt a minimal closure as the lowest-order local form consistent with locality, linearity in enclosed inertial content, rotational invariance, and additivity. These conditions restrict the sourced field to a divergence relation as the leading-order connection between field structure and inertial content, rather than an assumed inverse-square form. The concentration field therefore satisfies a divergence proportional to a local inertial density, with enclosed content given by its volume integral. In the weak-field limit, this density is dominated by rest mass and reduces to the baryonic density. All intrinsic kinematic scales are fixed by the acceleration scale $A$, while the constant $G$ sets the coupling to localized inertial content without introducing an additional kinematic scale.

For a spherically symmetric configuration the sourced acceleration is radial,
\begin{equation}
\mathbf a_{\mathrm{conc}}(r)=a_r(r)\,\hat r.
\end{equation}

In the local-reduction regime relevant for gravitationally bound systems, the concentration field resides in three spatial dimensions, so spherical symmetry is defined with respect to ordinary 2-spheres in $\mathbb{R}^3$. The flux through a sphere of radius $r$ is
\begin{equation}
\Phi(r)=\oint_S \mathbf a_{\mathrm{conc}}\cdot d\mathbf S
=4\pi r^2 a_r(r).
\end{equation}
Imposing flux conservation together with linearity in the enclosed inertial content gives
\begin{equation}
\oint_S \mathbf a_{\mathrm{conc}}\cdot d\mathbf S = -4\pi G M(r),
\label{eq:aft_flux_integral}
\end{equation}
where $M(r)$ is the enclosed inertial content within radius $r$, obtained from the local density by
\begin{equation}
M(r)=\int_V \rho_{\mathrm{inertial}}\,dV,
\end{equation}
and $G$ is the empirical coupling relating inertial content to concentration of the background acceleration. Spherical symmetry then yields
\begin{equation}
\mathbf a_{\mathrm{conc}}(r)=-\frac{G M(r)}{r^2}\,\hat r,
\label{eq:aft_inverse_square}
\end{equation}
and Gauss's theorem gives
\begin{equation}
\nabla\!\cdot \mathbf{a}_{\mathrm{conc}} = -4\pi G \rho_{\mathrm{inertial}},
\label{eq:aft_gauss_law}
\end{equation}
which is the Poisson equation for the sourced weak-field limit. In the weak-field regime relevant for galactic dynamics, the inertial density is dominated by baryonic rest-mass contributions, so $\rho_{\mathrm{inertial}}\approx\rho_b$.

\subsection{Circular Orbits}

For a body in a circular orbit of radius $r$ with tangential speed $v$, the required centripetal acceleration is
\begin{equation}
a_{\mathrm{cent}}(r) = -\frac{v^2}{r} .
\end{equation}
Equating this with the inward acceleration from equation~(\ref{eq:ain_scalar}) gives
\begin{equation}
-\frac{v^2}{r} = -A - \frac{G M_b(r)}{r^2} ,
\end{equation}
and multiplying by $r$ gives
\begin{equation}
v^2 = A r + \frac{G M_b(r)}{r} .
\end{equation}
Then, solving for the enclosed baryonic mass gives
\begin{equation}
M_b(r) = \frac{r\left(v^2 - A r\right)}{G}.
\label{eq:massradius}
\end{equation}

\subsection{Circular-Orbit Solutions}

Equation~(\ref{eq:massradius}) defines the PME fundamental plane, a two-parameter family of circular-orbit solutions relating baryonic mass $M$, orbital radius $r$, and orbital speed $v$.

For fixed orbital speed $v$, the mass function $M(r)$ is a concave quadratic in $r$. Differentiating gives
\begin{equation}
\frac{dM}{dr} = \frac{v^2 - 2 A r}{G} .
\end{equation}
The mass therefore reaches a maximum at
\begin{equation}
r_{\max} = \frac{v^2}{2A} .
\end{equation}
Substituting this radius into equation~(\ref{eq:massradius}) yields the maximum baryonic mass compatible with a circular orbit at speed $v$,
\begin{equation}
M_{\max}(v) = \frac{v^4}{4 A G}.
\label{eq:btfr_envelope}
\end{equation}

The surface $M(v,r)$ therefore defines a fundamental plane of circular-orbit solutions, as shown in Figure~\ref{fig:fundamentalplane}, with the ridge $M_{\max}(v)$ giving the upper envelope of baryonic mass for a given orbital speed. Galaxies dominated by circular rotation are expected to saturate this boundary and follow a characteristic mass--velocity relation.

\begin{figure}
\centering
\includegraphics[width=0.8\linewidth]{FundamentalPlane.pdf}
\caption{PME fundamental plane defined by equation~(\ref{eq:massradius}). The baryonic mass surface $M(v,r)$ is shown as a function of orbital speed $v$ and orbital radius $r$ for an illustrative manifold acceleration scale $A = 10^{-11}\,\mathrm{m\,s^{-2}}$. The red ridge marks the locus $M_{\max}(v)$ corresponding to the maximum baryonic mass permitted by the manifold kinematics at each orbital speed. Galaxies populating this ridge reproduce the baryonic Tully--Fisher relation.}
\label{fig:fundamentalplane}
\end{figure}

\subsection{Baryonic Tully--Fisher Relation}

$M_{\max}(v)$ defines the scaling $M \propto v^4$, consistent with the observed baryonic Tully--Fisher relation~\cite{McGaugh2000}. The manifold acceleration $A$ is determined empirically by fitting the circular-orbit envelope (equation~(\ref{eq:btfr_envelope})) to the SPARC galaxy sample of McGaugh, Lelli and Schombert~\cite{McGaugh2016}, using the assumptions $\Upsilon_\star=0.5$, $f_{\mathrm{gas}}=1.33$, and the quality cut $\mathrm{Qual}\leq 3$. These rotationally supported disk systems are well suited to the analysis, as their kinematics closely follow the circular-orbit approximation.

For each galaxy the rotation curve provides an asymptotic circular velocity $v$, while the baryonic mass $M_b$ is computed as $M_b=\Upsilon_\star L_{3.6}+f_{\mathrm{gas}}\,M_{\mathrm{HI}}$. The value of $A$ is determined by minimizing the $\chi^2$ statistic between the theoretical envelope and the observed $(M_b,v)$ pairs. The best-fit manifold acceleration is
\begin{equation}
A = (3.7 \pm 0.1)\times10^{-11}\,\mathrm{m\,s^{-2}} ,
\end{equation}
where the uncertainty is statistical, with $\Upsilon_\star$ and $f_{\mathrm{gas}}$ held fixed. The observed galaxies follow the predicted $M\propto v^4$ scaling across a wide range of baryonic masses and rotational velocities, as shown in Figure~\ref{fig:btfrchart}. A characteristic acceleration scale in galaxy dynamics has previously been emphasized in frameworks such as MOND~\cite{Milgrom1983}, in which the force law is modified phenomenologically to reproduce Tully--Fisher-type scaling. In PME, by contrast, the baryonic Tully--Fisher relation is a direct consequence of the constant-acceleration structure of the manifold and an observational signature of the acceleration flux.

\begin{figure}
\centering
\includegraphics[width=0.75\linewidth]{BTFRChart.pdf}
\caption{Log--log plot of baryonic mass $M_b$ versus outer rotation speed $v$ for the SPARC galaxy sample of McGaugh, Lelli and Schombert~\cite{McGaugh2016}. Black circles show the observed galaxies. The solid red line shows the PME circular-orbit envelope evaluated at the best-fit acceleration.}
\label{fig:btfrchart}
\end{figure}

\subsection{Global Acceleration Budget}

The homogeneous manifold carries a background acceleration field of magnitude $A$ at every point. Integrating this uniform field over the boundary of the homogeneous domain of scale $L$ gives a geometric flux fixed entirely by the manifold,
\begin{equation}
\Phi_{\mathrm{bg}} = -4\pi A L^2 .
\end{equation}
The local closure establishes that inertial content does not source new flux but redistributes the existing background into spatial concentrations. In the weak-field late-time regime, inertial content is dominated by baryonic rest mass, so $\rho_{\mathrm{inertial}} \approx \rho_b$. The total flux is therefore conserved: the sourced flux through the same boundary must equal the geometric flux of the empty manifold,
\begin{equation}
\Phi_{\mathrm{sourced}} = \oint \mathbf{a}_{\mathrm{conc}} \cdot d\mathbf{S} = -4\pi G M_b^{\mathrm{tot}} = -4\pi A L^2 .
\end{equation}
This conservation condition uniquely fixes the total baryonic mass within the homogeneous domain,
\begin{equation}
M_b^{\mathrm{tot}} = \frac{A L^2}{G} ,
\label{eq:total_baryonic_mass}
\end{equation}
as the mass for which the sourced and geometric fluxes are consistent. The baryonic content of the universe is therefore not a free parameter but a geometric consequence of the manifold acceleration scale $A$ and extent $L$. Local gravitational systems correspond to partial redistributions of this finite budget, so gravitational collapse approaches a saturated configuration rather than an unbounded one.

\subsection{Regularity of the Manifold}

The flux conservation law established in the previous subsection fixes the total acceleration flux available for redistribution at $-4\pi AL^2$. This is a finite geometric constant. Gravitational collapse concentrates flux but cannot create it: the sourced field is bounded above by the total budget, so collapse approaches a saturated configuration of finite density and finite radius. There is no point at which the field equations break down, no curvature divergence, and no geometric endpoint at which the evolution of the manifold terminates. Singularities are not hidden by a censorship mechanism — they do not arise.

Causal ordering is equally well-behaved. Separations are defined between ordered event pairs, and the operational clock accumulates monotonically along the manifold tangent. No trajectory can close back on itself without reversing this accumulation, so the framework admits no closed timelike curves. The causal structure is globally consistent at every epoch and at every scale.

The manifold is therefore smooth, causally ordered, and singularity-free at all times.

\section{Local Dynamics}
Having established the gravitational and causal structure, we now define the operational dynamics governing motion within the manifold. The treatment here establishes the kinematic foundation; further development of time-dilation phenomena and their geometric origin in the acceleration flux structure is outside the scope of this article.

\subsection{Operational Time}

The null condition established in equation~(\ref{eq:real_tangent_scale}) fixes the relation between spatial separation and temporal increment through the observable null scale $(V_0 - At)$,
\begin{equation}
\Delta t=\frac{\Delta x}{(V_0 - At)}.
\label{eq:operational_time}
\end{equation}
This defines the null boundary of the local geometry and identifies the observable null scale $(V_0-At)$ as the real projection of the underlying imaginary manifold tangent.

Observable time is defined operationally by null exchange: one tick is a complete signal round trip, and $t$ counts such cycles, defining an invariant ordering of events. Causality, motion, and measurable duration are thereby unified as real operational projections of the null structure of the manifold.

\subsection{Proper Time}
Using the local interval defined in equation~(\ref{eq:local_interval}),
\begin{equation}
ds^2 = -(V_0 - At)^2\,dt^2 + dx^2,
\end{equation}
proper time defines the invariant timelike interval along admissible trajectories, $d\tau^2 \equiv -ds^2$, so that
\begin{equation}
d\tau^2 = (V_0 - At)^2\,dt^2 - dx^2,
\end{equation}
giving
\begin{equation}
d\tau = (V_0 - At)\sqrt{1 - \frac{v^2}{(V_0 - At)^2}}\;dt,
\end{equation}
where $v = dx/dt$. Proper time therefore represents the timelike operational projection of the underlying imaginary tangent generator onto admissible observable trajectories.

\subsection{Local Momentum}

The homogeneous manifold acceleration $A$ selects a timelike direction, and the local interval (equation~(\ref{eq:local_interval})) decomposes the imaginary tangent evolution into orthogonal temporal and spatial components. Observable local geometry therefore resolves as real operational projections of this common tangent structure, with orientation parameterized by an angle $\theta$ relative to the flux.

Let $d\lambda$ denote an affine increment along the imaginary tangent, an invariant manifold parameter from which observable projections derive. Resolving the local projection structure induced by the tangent yields
\begin{equation}
dx = (V_0 - At)\sin\theta\,d\lambda,\qquad dt = \cos\theta\,d\lambda,
\end{equation}
with the observable null scale $(V_0-At)$ setting the projection scale per unit increment of $\lambda$ at epoch $t$. Measurable velocity is then the ratio of the spatial and temporal projections,
\begin{equation}
v = \frac{dx}{dt} = (V_0 - At)\tan\theta,
\end{equation}
and observable momentum follows by inertial weighting,
\begin{equation}
P = m\,v = P_{\mathrm{geom}}\tan\theta,
\end{equation}
where $P_{\mathrm{geom}} \equiv m(V_0-At)$ is the geometric momentum scale associated with the observable null projection of the manifold tangent at epoch $t$.

The null boundary is reached when $v = V_0 - At$, requiring $\tan\theta = 1$ and therefore $\theta = \pi/4$. At this orientation the spatial and temporal projections are equal in magnitude, and observable motion is restricted to the timelike domain $0 \le \theta < \pi/4$. Beyond this boundary the projection geometry remains well-defined through $\theta$, but the kinematic description in terms of $t$ ceases to apply; its continuation is defined in terms of phase evolution.

\subsection{Local Action}

The bilocal construction admits a natural system of units in which mass and energy are reciprocally dimensioned along the manifold tangent. Mass is rendered in units of time per unit length, while energy comes in units of length per unit time. Under this convention, the product $m(V_0-At)$ is dimensionless and energy has units of velocity, same as the manifold tangent. The dimensional conversion constants between inertial and recurrence quantities are therefore absorbed into the geometry rather than introduced separately.

In the coincidence limit, the local interval (equation~(\ref{eq:local_interval})) defines the observable projection structure induced by the underlying imaginary tangent evolution together with an invariant interval $ds$ describing local manifold separation. Weighting this invariant separation by the geometric momentum scale
\begin{equation}
P_{\mathrm{geom}} \equiv m(V_0-At),
\end{equation}
defines the invariant geometric action increment,
\begin{equation}
dS_{\mathrm{geom}} \equiv P_{\mathrm{geom}}\,ds,
\end{equation}
with quadratic form
\begin{equation}
dS_{\mathrm{geom}}^2 = P_{\mathrm{geom}}^2\,ds^2.
\end{equation}

Using the local invariant
\begin{equation}
ds^2 = - (V_0 - At)^2 dt^2 + dx^2,
\end{equation}
the action-space interval becomes
\begin{equation}
dS_{\mathrm{geom}}^2 =m^2(V_0-At)^2\left[-(V_0-At)^2 dt^2 + dx^2\right].
\end{equation}

The temporal component of $dS_{\mathrm{geom}}$ carries a factor of $-(V_0-At)^2 dt^2$, so the action rate per unit $dt$ in the temporal sector is $-m(V_0-At)^2$. This defines the geometric energy scale,
\begin{equation}
E_{\mathrm{geom}}(t)\equiv -m(V_0-At)^2,
\label{eq:geometric_energy_definition}
\end{equation}
while the spatial component is governed by the same geometric momentum scale $P_{\mathrm{geom}}=m(V_0-At)$.

Evaluating the local interval along the projected tangent trajectory gives $ds^2 = -(V_0-At)^2 d\lambda^2$, which is timelike on the expansion branch. Taking the positive root yields the invariant element along the tangent,
\begin{equation}
ds=(V_0-At)\,d\lambda,
\end{equation}
and therefore
\begin{equation}
dS_{\mathrm{geom}} = m(V_0-At)^2\,d\lambda,
\label{eq:action_definition}
\end{equation}
so the invariant geometric action accumulates along the projected manifold trajectory,
\begin{equation}
S_{\mathrm{geom}} = \int dS_{\mathrm{geom}}.
\end{equation}

Under the geometric unit convention established above, $dS_{\mathrm{geom}}$ inherits the dimensions of invariant separation itself. Action is therefore expressed geometrically in units of length, so any minimum action scale corresponds geometrically to a minimum manifold separation rather than an independently imposed dimensional constant.

The invariant geometric action $dS_{\mathrm{geom}}$ admits a projection into the observable timelike sector, defining an observable action increment $dS$. Resolving the tangent evolution into its temporal and spatial contributions expresses this projection in canonical form,
\begin{equation}
dS = P\,dx - H\,dt,
\end{equation}
which identifies $P$ as the observable momentum conjugate to $x$ and $H$ as the observable Hamiltonian conjugate to $t$. The Hamiltonian $H$ coincides with the geometric energy scale $E_{\mathrm{geom}}$ defined in equation~(\ref{eq:geometric_energy_definition}). The observable action $dS$ therefore admits a Hamilton--Jacobi interpretation within the timelike projection sector, while $dS_{\mathrm{geom}}$ remains the invariant quantity defined across the full manifold.

Since $dS_{\mathrm{geom}}$ defines the momentum-weighted invariant separation of the manifold, physical trajectories follow the shortest paths in action space. The principle of least action therefore emerges as the geometric shortest-distance condition on the invariant action geometry.

\subsection{Operational Energy}
The operational clock counts ordered null exchanges between neighboring worldlines, with one tick corresponding to a complete signal round trip. For an event pair separated by $\Delta x$, the operational clock definition of equation~(\ref{eq:operational_time}) gives a clock frequency
\begin{equation}
\nu = \frac{1}{\Delta t} = \frac{V_0-At}{\Delta x},
\end{equation}
and multiplying by the separation recovers the invariant tangent scale,
\begin{equation}
\Delta x\,\nu = V_0-At.
\end{equation}
Because the bilocal geometry admits no infinitesimal separations as fundamental geometric objects, operational frequency possesses physical meaning only across a minimum bilocal separation, denoted $h$. Evaluating the null exchange at this minimum scale defines the operational energy,
\begin{equation}
E_{\mathrm{op}} \equiv h\nu,
\label{eq:operational_energy_relation}
\end{equation}
which is the Planck--Einstein relation, recovered here as a geometric consequence of the bilocal construction rather than a postulated quantization. At the operational floor, $\nu = (V_0-At)/h$, so the operational energy reduces to the tangent invariant,
\begin{equation}
E_{\mathrm{op}} = h\nu = V_0 - At,
\label{eq:operational_floor_identity}
\end{equation}
identifying $h$ not as an independently introduced constant but as a geometric property of the bilocal construction: the minimum separation at which the operational clock structure coincides directly with the underlying manifold tangent. Under the geometric unit convention, $h$ and the invariant action $S_{\mathrm{geom}}$ both carry the dimensions of geometric separation, so the Planck scale arises as a geometric length rather than an independent dimensional input.

The geometric energy $E_{\mathrm{geom}}(t) \equiv -m(V_0-At)^2$ is the inertially weighted counterpart of the same tangent structure. The two are related by the geometric momentum scale $P_{\mathrm{geom}} = m(V_0-At)$,
\begin{equation}
E_{\mathrm{geom}} = -P_{\mathrm{geom}}\,E_{\mathrm{op}},
\label{eq:energy_unification}
\end{equation}
so geometric and operational energy arise from a common invariant tangent structure, differing only in inertial weighting. Geometric energy is the top-down projection of the manifold onto inertial trajectories; operational energy is the bottom-up counting structure of the bilocal clock. Equation~(\ref{eq:energy_unification}) is their unification.

\subsection{Phase Evolution}
The projection-based kinematic description applies within the timelike sector, where the local invariant defines a real proper-time interval. At the null boundary, $ds^2=0$ and $d\tau \to 0$, marking the limit of observable timelike motion. Beyond this boundary, observable motion ceases, while the invariant geometric structure remains well-defined and the geometric action continues to accumulate along the manifold tangent.

Beyond the null boundary, the underlying imaginary tangent generator persists even though observable timelike parametrization ceases within the real operational sector. Phase therefore parametrizes the continued accumulation of invariant geometric action beyond the proper-time sector. Because the operational clock is defined through ordered null recurrence, one complete operational bounce subtends $2\pi$ radians of accumulated phase. The corresponding action-per-radian conversion is therefore
\begin{equation}
\hbar \equiv \frac{h}{2\pi},
\end{equation}
where $h$ is the minimum operationally meaningful bilocal separation defined by the operational floor identity~(\ref{eq:operational_floor_identity}). Phase evolution is then given by
\begin{equation}
\phi = \frac{S_{\mathrm{geom}}}{\hbar}.
\end{equation}
Combined with equation~(\ref{eq:operational_floor_identity}), this gives $\hbar\omega = V_0 - At$ at the operational floor, with $\omega = 2\pi\nu$, recovering the angular form of the Planck--Einstein relation as a geometric closure condition on the operational cycle rather than as a postulated quantization. The distinction between $h$ and $\hbar$ is physical rather than notational in the present framework: $h$ defines the minimum operational separation underlying the recurrence structure of the manifold, while $\hbar$ converts accumulated geometric action into radian-valued phase through closure of the operational cycle.

Phase evolution therefore represents the continuation of invariant manifold evolution beyond the proper-time-parametrized sector, while the underlying geometric structure remains continuous across the null boundary.

\section{Late-Time Evolution}

The distance--redshift relation derived from the bilocal interval predicts the late-time expansion history.

\subsection{Type Ia Supernova Sample}

We use the Pantheon+ compilation of 1701 spectroscopically confirmed Type~Ia supernovae presented by Brout et~al.~\cite{Brout2022}. The publicly released Pantheon+SH0ES dataset provides the observed distance moduli $\mu_{\mathrm{obs}}$, redshifts $z$, and the full STAT+SYS covariance matrix for covariance-weighted $\chi^2$ inference. The distance moduli are absolutely calibrated through the SH0ES determination of the Type~Ia absolute magnitude $M_B$.

\subsection{Luminosity Distance}

The spatial separation predicted by the manifold geometry is given by equation~(\ref{eq:distance_redshift_relation}). The luminosity distance follows from the standard relation $D_L(z) = (1+z)D_C(z)$. Substituting the spatial separation yields the corresponding analytic luminosity distance
\begin{equation}
D_L(z) =-\frac{t_o\left(A t_o - 2V_0\right) z(1+z)}{2+z}.
\label{eq:pme_luminosity_distance}
\end{equation}

\subsection{Distance Modulus}

Type~Ia supernova observations are reported in terms of the distance modulus, which relates the observed luminosity distance to the measured brightness of each supernova,
\begin{equation}
\mu(z) = 5\log_{10}\!\left(\frac{D_L(z)}{10\,\mathrm{pc}}\right).
\end{equation}

\subsection{Manifold Kinematic Parameters}

The local null scale at the observation epoch is
\begin{equation}
c_o = V_0 - A t_o .
\label{eq:co_definition}
\end{equation}
Operationally, photon exchange is used as a proxy for this local null scale, so the measured speed of light at the observation epoch determines $c_o$, fixing the integration constant,
\begin{equation}
V_0 = c_o + A t_o .
\label{eq:velocity_boundary_condition}
\end{equation}

With the manifold acceleration $A$ determined from the baryonic Tully--Fisher analysis, the Cepheid-calibrated Pantheon+ subset is used to determine the observation epoch $t_o$ on the absolute distance scale. For each candidate value of $t_o$, the corresponding $V_0$ follows algebraically from equation~(\ref{eq:velocity_boundary_condition}), and theoretical distance moduli $\mu_{\mathrm{th}}(z)$ are evaluated at the calibration-set redshifts. These are compared with the observed values $\mu_{\mathrm{obs}}(z)$, defining the residual vector
\begin{equation}
r_i = \mu_{\mathrm{obs}}(z_i) - \mu_{\mathrm{th}}(z_i).
\end{equation}
The best-fit observation epoch is obtained by minimizing the covariance-weighted chi-square statistic
\begin{equation}
\chi^2 = r^{\mathrm T} C^{-1} r,
\end{equation}
where $C$ is the full Pantheon+ STAT+SYS covariance matrix. The resulting best-fit value of $t_o$ then determines $V_0$, completing the kinematic parameter set. The resulting parameters are listed in Table~\ref{tab:pme_parameters}.

\begin{table}
\centering
\begin{tabular}{|l|l|}
\hline
Parameter & Value \\
\hline
$A$ & $3.7\times10^{-11}\,\mathrm{m\,s^{-2}}$ \\
$V_0$ & $3.16\times10^{8}\,\mathrm{m\,s^{-1}}$ \\
$t_o$ & $4.51\times10^{17}\,\mathrm{s}$ \\
\hline
\end{tabular}
\caption{PME kinematic parameters determined from BTFR and Cepheid-calibrated Pantheon+ subset.}
\label{tab:pme_parameters}
\end{table}

\subsection{Comparison with Observations}

\begin{figure}
\centering
\includegraphics[width=0.7\linewidth]{HubbleDiagram.pdf}
\caption{Pantheon+SH0ES distance moduli as a function of redshift $z$, compared with the PME prediction and the spatially flat FLRW model. PME parameters are fixed by BTFR and Cepheid calibration. FLRW parameters are fixed by the Planck 2018 results.}
\label{fig:hubblechart}
\end{figure}

Figure~\ref{fig:hubblechart} shows the PME prediction and the Planck 2018 FLRW model compared with the Pantheon+SH0ES data.

The Pantheon+SH0ES dataset provides observed distance moduli $\mu_{\mathrm{obs}}(z)$ calibrated on an absolute scale through the SH0ES determination of the Type~Ia supernova absolute magnitude. This calibration is applied identically to both models.

Each model predicts $D_L(z)$, which is mapped to $\mu_{\mathrm{th}}(z)$ using the definition of $\mu(z)$ introduced earlier. We compare $\mu_{\mathrm{th}}(z)$ to $\mu_{\mathrm{obs}}(z)$ using the full Pantheon+ STAT+SYS covariance matrix.

PME is evaluated using the fixed parameter set derived above. FLRW is evaluated using a spatially flat model specified by the Planck 2018 parameter set $\Omega_m = 0.315$, $\Omega_\Lambda = 1-\Omega_m$, and $H_0 = 67.4\,\mathrm{km\,s^{-1}\,Mpc^{-1}}$~\cite{Planck2018}. The comparison is performed on the out-of-sample supernovae, excluding the Cepheid calibration records.

The resulting covariance-weighted chi-square values are
\begin{equation}
\chi^2_{\mathrm{PME}}=2726,\qquad \chi^2_{\mathrm{FLRW}}=4049,
\end{equation}
for $\nu=1625$, corresponding to reduced values $\chi^2/\nu=1.68$ and $2.49$, respectively. PME predicts the observed distance moduli substantially more accurately than the Planck 2018 FLRW model across the Pantheon+SH0ES sample, without dark sector parameters.

\subsection{Hubble Function}

The Hubble expansion rate is defined by the fractional rate of change of the real spatial extent of the manifold, $H(t) \equiv \dot{L}(t)/{L(t)}$. Differentiating equation~(\ref{eq:manifold_extent}) gives $\dot{L}(t)=V_0-At$, so
\begin{equation}
H(t) = \frac{V_0 - At}{V_0 t - \frac{1}{2}A t^2}.
\label{eq:hubble_function}
\end{equation}

Evaluating equation~(\ref{eq:hubble_function}) at the observation epoch $t=t_o$ using the best-fit manifold parameters gives the present-day expansion rate
\begin{equation}
H_0 = 66.5\,\mathrm{km\,s^{-1}\,Mpc^{-1}}.
\end{equation}
The Hubble constant therefore emerges as a prediction rather than a fitted parameter.

\subsection{Cosmic Time}

The redshift–time relation differs from standard cosmological evolution due to the underlying parabolic evolution. Table~\ref{tab:cosmic_time_comparison} compares the resulting cosmic times with those from FLRW algorithms.

\begin{table}
\centering
\begin{tabular}{|c|c|c|}
\hline
$z$ & PME [Myr] & FLRW [Myr] \\
\hline
0    & 14,283 & 13,791 \\
1    & 7,039 & 5,840 \\
10   & 1,265 & 470 \\
100  & 137 & 16.4 \\
1,000 & 13.9 & 0.429 \\
\hline
\end{tabular}
\caption{Cosmic time as a function of redshift. PME predictions are compared with FLRW values computed using Planck 2018 parameters~\cite{Planck2018}.}
\label{tab:cosmic_time_comparison}
\end{table}

The parabolic expansion yields uniformly larger cosmic ages than FLRW, significantly extending the available time for early structure formation.

\section{Early-Time Evolution}

Independent constraints arise from the cosmic microwave background (CMB), whose acoustic scale probes the distance to last scattering relative to the sound horizon.

The expansion history is determined by the PME kinematics derived above, while recombination microphysics is evaluated using standard rate equations. Thermodynamic quantities are mapped using the scale factor
\begin{equation}
a(t) = \frac{L(t)}{L(t_o)} .
\label{eq:scale_factor}
\end{equation}

\subsection{Baryon and Photon Densities}

The recombination history and pre-decoupling sound speed require the baryon and photon densities. The baryon density is fixed by the manifold acceleration scale rather than introduced as an independent cosmological parameter. This follows from the local divergence law (equation~(\ref{eq:aft_gauss_law})), which relates the sourced acceleration field to the inertial density. In the weak-field late-time regime relevant here, the inertial density is dominated by baryonic rest mass, $\rho_{\mathrm{inertial}} \approx \rho_b$, so the baryon density is determined directly from the manifold parameters.

For a homogeneous baryonic distribution of density $\rho_b$ within a region of characteristic size $L$, the enclosed mass is
\begin{equation}
M(L)=\frac{4\pi}{3}L^3\rho_b ,
\end{equation}
so the corresponding baryon-sourced acceleration at radius $L$ is
\begin{equation}
\frac{GM(L)}{L^2}=\frac{4\pi}{3}G L\rho_b .
\end{equation}
Evaluating this expression at the present manifold extent $L(t_o)$, which contains the full homogeneous baryonic source, and identifying it with the inferred manifold acceleration $A$ gives
\begin{equation}
A=\frac{4\pi}{3}G L(t_o)\rho_b ,
\end{equation}
hence
\begin{equation}
\rho_b=\frac{3A}{4\pi G\,t_o\left(V_0-\frac{1}{2}A t_o\right)} .
\end{equation}

Using the kinematic parameters determined in the previous sections yields
\begin{equation}
\rho_b = 9.49\times10^{-28}\,\mathrm{kg\,m^{-3}}.
\end{equation}
This value is consistent with independent observational estimates of the present-day baryon density from primordial light-element abundances~\cite{Cooke2018}.

The photon density is determined from the measured CMB temperature $T_\gamma=2.725\,\mathrm{K}$. The mapping between temperature and photon density depends on the local geometric scale $(V_0 - At)$, so the photon density is treated as epoch dependent rather than constant. For a blackbody radiation field, the photon number density is
\begin{equation}
n_\gamma(t)=\frac{2\zeta(3)}{\pi^2}\left(\frac{k_B T_\gamma}{\hbar (V_0-At)}\right)^3 ,
\end{equation}
and, following the geometric sign convention of Eq.~(\ref{eq:geometric_energy_definition}), the mean photon energy is
\begin{equation}
\langle E_\gamma\rangle=-\frac{\pi^4}{30\zeta(3)}k_B T_\gamma .
\end{equation}
The corresponding photon energy density is
\begin{equation}
u_\gamma(t)=n_\gamma(t)\langle E_\gamma\rangle ,
\end{equation}
so the mass-equivalent photon density becomes
\begin{equation}
\rho_\gamma(t)=\frac{u_\gamma(t)}{-(V_0-At)^2}.
\end{equation}

These densities determine the baryon--photon momentum ratio and the free-electron density entering the visibility function.

\subsection{Recombination Epoch}

The recombination epoch is the time of maximum photon visibility, i.e. the most probable last-scattering time in the baryon--photon plasma.

To determine the ionization history, we evaluate the standard recombination rate equations using \texttt{Recfast++} with the PME expansion geometry. The background evolution is provided by the PME Hubble function (equation~(\ref{eq:hubble_function})) in place of the FLRW algorithm, with the expansion rate governed by the best-fit manifold parameters. The baryon sector is specified by the present-day baryon density derived above. The resulting ionization history is used to interpolate the free-electron fraction $x_e(z)$ entering the visibility function. The free-electron number density $n_e(t)$ is defined as
\begin{equation}
n_e(t)=x_e(t)\,\frac{\rho_b}{m_p}\,a(t)^{-3},
\end{equation}
where $\rho_b$ is the present-day baryon mass density, $m_p$ is the proton mass.

The Thomson scattering rate is expressed in terms of the local particle density, cross section, and local geometric scale,
\begin{equation}
\Gamma_T(t)=n_e(t)\,\sigma_T\,(V_0 - At),
\end{equation}
where $\sigma_T$ is the Thomson cross section. The optical depth is then defined as the accumulated interaction probability along the manifold history,
\begin{equation}
\tau_T(t)=\int_t^{t_o}\Gamma_T(t')\,dt'.
\end{equation}
The corresponding visibility function is
\begin{equation}
g(t)=\Gamma_T(t)\,e^{-\tau_T(t)} .
\end{equation}

Evaluating $g(t)$ on the PME expansion background using the \texttt{Recfast++} microphysics parametrization yields a maximum at $z_* = 1039$, $t_* = 13.4\,\mathrm{Myr}$ shown in Fig.~\ref{fig:visibility}.

\begin{figure}
\centering
\includegraphics[width=0.75\linewidth]{VisibilityFunction.pdf}
\caption{Photon visibility function $g(z)$ evaluated using the recombination history on the PME expansion background. The peak defines the recombination epoch, yielding $z_* = 1039$.}
\label{fig:visibility}
\end{figure}

\subsection{Sound Horizon}

Prior to recombination, baryons and photons form a tightly coupled fluid whose acoustic interaction distance defines the sound horizon. The thermodynamic evolution follows the scale factor $a(t)$ defined in equation~(\ref{eq:scale_factor}).

The baryon--photon momentum ratio is
\begin{equation}
R(t)=\frac{3\rho_b}{4\rho_\gamma(t)}\,\frac{a(t)^{-3}}{a(t)^{-4}}=\frac{3\rho_b}{4\rho_\gamma(t)}\,a(t),
\end{equation}
giving the sound speed
\begin{equation}
c_s(t)=\frac{c_{\mathrm{c}}(t)}{\sqrt{3\left(1+R(t)\right)}} .
\end{equation}
The causal expansion rate $c_{\mathrm{c}}(t)$ sets the rate at which interactions can be ordered through the operational clock and thus defines the effective causal scale entering the acoustic horizon integral, so that the sound horizon is
\begin{equation}
r_s=\int_0^{t_*} c_s(t)\,dt .
\end{equation}
Evaluating this expression using the manifold parameters yields $r_s=3.7\,\mathrm{Mpc}$.

\subsection{Acoustic Angular Scale}

The observed acoustic angle is determined by the ratio of the sound horizon at recombination to the bilocally assigned transverse curvature scale of the emission--observation pair. Because the spatial slices are closed with circumference $L(t)$, the acoustic feature at recombination spans the full transverse extent of the last-scattering slice, and the observer integrates over the full transverse extent of the present slice. The bilocal pair relevant to this observation is therefore the maximally-bracketed configuration in which the transverse endpoints on each slice reach the full slice circumference, so that $\Delta x_e^1 = L(t_*)$ and $\Delta x_o^1 = L(t_o)$.

Applying the bilocal spatial rule of equation~(\ref{eq:spatial_half_expansion}) to this maximally-bracketed transverse configuration defines the bilocal-averaged transverse circumference,
\begin{equation}
C_\perp \equiv \frac12\left[L(t_o)+L(t_*)\right].
\label{eq:bilocal_transverse_circumference}
\end{equation}
Using the redshift relation from equation~(\ref{eq:redshift_relation}), this becomes
\begin{equation}
C_\perp=\frac{L(t_o)}{2}\,\frac{2+z_*}{1+z_*}.
\end{equation}

Angular features on a closed transverse surface subtend angles relative to its radius of curvature. For the bilocal pair the effective transverse radius is therefore obtained from the bilocal-averaged circumference,
\begin{equation}
r_\perp=\frac{C_\perp}{2\pi}=\frac{L(t_o)}{4\pi}\,\frac{2+z_*}{1+z_*}.
\end{equation}

The acoustic angular scale is then
\begin{equation}
\theta_*=\frac{r_s}{r_\perp}.
\end{equation}
Substituting the expression above for $r_\perp$ gives
\begin{equation}
\theta_*=\frac{4\pi(1+z_*)}{L(t_o)(2+z_*)}\,r_s.
\end{equation}
Using equation~(\ref{eq:manifold_extent}) evaluated at $t_o$ yields
\begin{equation}
\theta_*=\frac{8\pi(1+z_*)}{t_o(2V_0-A t_o)(2+z_*)}\,r_s.
\label{eq:acoustic_angle_final}
\end{equation}

Evaluating this expression with the manifold parameters gives
\begin{equation}
\theta_* = 0.0103.
\end{equation}
Expressed in the conventional CMB notation,
\begin{equation}
100\,\theta_* = 1.03.
\end{equation}

This result follows from the PME expansion geometry combined with standard recombination microphysics.

\section{Discussion}

The minimum bilocal separation $h$ introduced in the construction of geometric energy is not an independently imposed discreteness postulate, but follows from the operational definition of the geometry itself. Because geometry is defined only through ordered bilocal relations, the framework contains no independently defined sub-operational continuum in which infinitesimal separations possess autonomous geometric meaning. The local continuum appears only as the coincidence-limit approximation of fundamentally finite bilocal relations. The divergence of the recurrence density in the coincidence limit therefore reflects the breakdown of the operational description as the bilocal separation ceases to remain physically meaningful. The bilocal geometry admits a minimum operational separation, represented by $h$, at which the operational energy coincides exactly with the manifold tangent scale $(V_0 - At)$, uniquely fixing $h$ as a geometric property rather than an independent dimensional input.

The momentum-weighted invariant separation defines an action geometry in which physical trajectories follow shortest paths in action space. The principle of least action is therefore recovered as a geometric shortest-distance condition rather than imposed as an independent variational postulate.

The time-dependent PME energy scale breaks the time-translation invariance implicit in standard BBN calculations and relaxes the equilibrium assumptions underlying baryon asymmetry arguments; a self-consistent early-universe treatment is outside the scope of this article. Extending the framework to perturbations, structure growth, and a fully intrinsic recombination treatment is similarly deferred.

The framework predicts a residual non-Keplerian contribution to bound motion with radial scaling $A r^2/GM$. In compact systems this effect is suppressed and becomes significant only at larger radii. Solar-System observations therefore constrain the model but do not probe the regime in which the background term becomes dynamically dominant. Improved precision, particularly at large heliocentric distances, should reveal systematic departures from purely Keplerian motion. Detection or exclusion of such a signal provides a direct test of the construction.

The three forms of energy identified in the framework — geometric, operational, and kinetic — are not independent postulates connected empirically, as in standard physics, but distinct readings of a single geometric object. Geometric energy $E_{\mathrm{geom}} = -m(V_0-At)^2$ is the inertially weighted temporal projection of the manifold tangent, the rest energy of the trajectory. Operational energy $E_{\mathrm{op}} = h\nu$ is the frequency of null exchange at the minimum bilocal separation, recovering the Planck--Einstein relation geometrically. Kinetic energy arises as the $\theta$-dependent projection of the same tangent onto spatial trajectories. Equation~(\ref{eq:energy_unification}) is the statement that all three are the same invariant tangent structure viewed from different operational perspectives.

\section{Code Availability}

A computational notebook implementing the bilocal geometry, distance--redshift relation, and observational analysis is publicly available at:
\begin{center}
\url{https://www.wolframcloud.com/obj/32b2e831-2cbe-4ff5-b821-aa23f476f015}
\end{center}
The notebook reproduces the results and figures presented in this work and is provided to enable independent verification of the calculations.

\section{Conclusion}

A constrained bilocal construction on a complex manifold yields a quadratic interval whose local reduction produces a homogeneous acceleration scale. This single geometric scale governs both the parabolic evolution of the manifold and the redistribution of acceleration flux in gravitationally bound systems. The bilocal geometry produces an emergent local null constraint, defines an invariant action geometry, and admits a real operational projection structure from which observable time, proper time, momentum, energy, and phase arise.

In the local-reduction regime, action is not an independently postulated functional but the momentum-weighted invariant separation of the manifold. Physical trajectories therefore follow shortest paths in action space, so the principle of least action emerges as a geometric consequence of the bilocal construction. Phase evolution provides the continuation of invariant tangent evolution beyond the proper-time-parametrized sector, while the operational recurrence structure identifies the role of $h$ and $\hbar$ within the geometry rather than introducing them as independent dimensional inputs.

The evolution implied by this construction yields a closed analytic distance--redshift relation, from which the Pantheon+ Type~Ia supernova luminosity distances follow directly on an absolute scale. The same parameter set predicts the present-day Hubble constant and produces uniformly extended cosmic ages. In the local-reduction regime, redistribution of the homogeneous acceleration field generates an inverse-square sourced component and a circular-orbit envelope consistent with the observed baryonic Tully--Fisher relation. The acceleration scale inferred from galaxy dynamics then fixes the baryonic density of the homogeneous universe, linking galactic and cosmological observables without introducing non-baryonic dark matter or dark energy components.

The bilocal construction also determines the causal horizon and implies a constant ratio between the causal domain and manifold extent, placing the last-scattering surface within a single causal region and accounting for large-scale isotropy. The geometry yields a transverse scale for recombination that reproduces the observed acoustic angular scale. Because the available acceleration flux is finite, gravitational collapse approaches saturated configurations, eliminating curvature singularities and enforcing a globally ordered causal structure. Rotational inertia emerges from redistribution of the same finite acceleration budget, providing a Machian interpretation in which local inertial behavior depends on the global distribution of inertial content.

Evolution, causal structure, action, gravity, inertia, recurrence, and phase therefore arise from a common bilocal origin, with the framework's local kinematic and gravitational structure arising from the momentum-weighted invariant separation. The framework provides a unified explanation of isotropy, galaxy rotation, Type~Ia supernova magnitudes, the acoustic scale, nonsingular gravitational collapse, the baryonic density, and the Hubble constant, while remaining directly testable through galactic dynamics, high-redshift distance measurements, gravitational lensing, Solar-System residuals, and the requirement that a complete early-universe treatment reproduce observed nucleosynthesis.

\bibliographystyle{unsrt}
\bibliography{references}

\end{document}